\section{Discussion}
\label{sec:discussion}

RQ1 has no scalar answer.  RNDVoC's range-null
decomposition produced the strongest intrusive-proxy state at modest
parameter count, Vocos converted frame-level spectral prediction into the
lowest RTF on both lanes, and RFWave traded iterative solver work for an
intermediate point; the small HiFi-GAN variants minimized footprint below
the leading quality band.  Generation granularity and analytical output
structure tracked the observed orderings better than parameter count
alone, a mechanism-consistent reading, not causal identification, under
native budgets and single trajectories.  The choice of quality proxy is itself a deployment decision: PESQ
and STOI selected RNDVoC while UTMOS selected HiFi-GAN V1 and reshaped
both cost frontiers; proxy agreement still would not constitute listener
preference.  RQ3 inherited this conditionality: the selected artifact
changed with lane and constraint vector.

RQ2 resolved conditionally on operator graph, backend, and
runtime.  FP16 casting halved storage almost uniformly, consistent with its
arithmetic-format origin \citep{micikevicius2018mixedprecision}, yet its
latency effect spanned $0.554\times$ to $1.955\times$; compilation and ONNX
gains depended on graph lowering and kernel choice
\citep{ansel2024pytorch2}; and INT8 coverage was necessary but
insufficient: transformed paths could shrink bytes yet regress speed
\citep{jacob2018quantization}.  Static ONNX INT8 also exposed calibration
dependence: execution-seeded recalibration varied the quantized artifact
by 0.14 PESQ for HiFi-GAN V1; deployed integer pipelines should treat
the calibration set as part of the artifact's identity.  With hosts unpaired, large contrasts are more
credible than marginal crossings, and the taxonomy remains an exploratory
summary of continuous effects.

Dense masks and solver reduction manipulated different objects.  Masks
measured tolerance to zeroed weights, not sparse deployment
\citep{han2015pruning}; recovery showed quality is relearnable,
BigVGAN-base regaining 1.17 PESQ, but the missing dense-continuation
control confounds recovery with added training.  RFWave's step ladder
changed the iteration count and produced the matrix's only smooth,
monotone speed-quality curve, making eight steps a defensible choice.  VocosFormer supports a bounded mechanism hypothesis:
unmasked global attention can spread padded-frame contamination, matching
the padded-to-true-length gap; true-length evaluation removed most of
the quality gap, none of the latency penalty, and attribution requires
masked- and local-attention controls.

In this cohort, post hoc compression split sharply: FP16 held PESQ at
half size in eight of ten groups, while dense masking harmed 18 of 19
groups and static ONNX INT8 three of four; the larger efficiency
separations were architectural, consistent with the architecture-first
emphasis of recent designs
\citep{siuzdak2024vocos,liu2025rfwave,lee2025periodwave,oh2026comvo}.

All intervals are fixed-set descriptive quantities conditional on the
adaptive single-speaker suffix, not an untouched test set; no
multiplicity correction, training replication, listening study, or
external evaluation exists, and warm timing omits streaming, queuing,
cold start, and host identity.  The priority follow-up is an untouched
multi-speaker test set with blinded listening and replicated
trajectories, then randomized host-paired timing with retained raw
repetitions, with sparse kernels, quantization-aware training,
pinned-calibration INT8, and attention masking in those controls.
